\documentclass[11pt]{article}

\usepackage[margin=1in]{geometry}
\usepackage{amsmath,amssymb}
\usepackage{array}
\usepackage{booktabs}
\usepackage{longtable}
\usepackage{hyperref}

\title{GraphFEX Complex-Network FEX Implementation Audit}
\author{xingjian\_check branch}
\date{September 10, 2026}

\begin{document}
\maketitle

\section{Purpose}

This note summarizes the implementation issues found in the GraphFEX
complex-network FEX path before running the proposed HR tests.  The goal is to
separate three things:
\begin{enumerate}
    \item whether the repository matches the published FEX algorithm,
    \item whether stratified scoring is implemented as a real controller
    scoring rule rather than a fixed-design proxy, and
    \item which issues must still be tested before making a paper claim.
\end{enumerate}

The audit used the current repository on branch \verb|xingjian_check| and the
online FEX references:
\begin{itemize}
    \item Network FEX: \href{https://arxiv.org/abs/2401.03092}{arXiv:2401.03092}.
    \item H-FEX: \href{https://arxiv.org/abs/2506.20607}{arXiv:2506.20607}.
    \item FEX/TFDM stochastic residual framework:
    \href{https://arxiv.org/abs/2605.10687}{arXiv:2605.10687}.
\end{itemize}

\section{Reference Objective}

For network dynamics, the target structure is
\begin{equation}
    \dot{x}_i(t)
    =
    F(x_i(t))
    +
    \sum_{j=1}^{N} A_{ij}G(x_i(t),x_j(t)).
\end{equation}
For a candidate operator sequence $e$ with trainable parameters $\theta$, the
full candidate loss is
\begin{equation}
    L_{\mathrm{full}}(e,\theta)
    =
    \frac{1}{T N}
    \sum_{t=1}^{T}\sum_{i=1}^{N}
    \left[
    \widehat{\dot{x}}_i(t)
    -
    F_{\theta}(x_i(t))
    -
    \sum_{j:(i,j)\in E} A_{ij}
    G_{\theta}(x_i(t),x_j(t))
    \right]^2.
\end{equation}
The controller score should be a monotone transformation of the optimized loss,
\begin{equation}
    S(e)=\frac{1}{1+L(e,\widehat{\theta}_e)}.
\end{equation}
The controller is then updated from the top quantile of sampled operator
sequences, and the stored candidates are fine-tuned after controller search.

\section{Scoring Variants}

\subsection{Random Batch Method}

For a random partition of nodes into batches, RBM keeps only edges whose two
endpoints fall in the same batch:
\begin{equation}
    (i,j)\in E_{\mathrm{RBM}}
    \quad \Longleftrightarrow \quad
    b(i)=b(j).
\end{equation}
If all batches have size $p$, a fixed edge is retained with probability
\begin{equation}
    q=\frac{p-1}{N-1}.
\end{equation}
Raw RBM is therefore a biased interaction estimator:
\begin{equation}
    \mathbb{E}\left[\widetilde{H}_i(t)\right]
    =
    q H_i(t),
    \qquad
    H_i(t)=\sum_j A_{ij}G(x_i(t),x_j(t)).
\end{equation}
The optional inverse-probability correction is
\begin{equation}
    A_{ij}\mapsto \frac{A_{ij}}{q}.
\end{equation}
For unequal batches of size $s_b$, the correct retained-edge probability is
\begin{equation}
    q_b=\frac{s_b-1}{N-1}.
\end{equation}

\subsection{Stratified Scoring}

For a selected node stratum $S$, the stratified controller loss is
\begin{equation}
    L_S(e,\theta)
    =
    \frac{1}{T|S|}
    \sum_{t=1}^{T}\sum_{i\in S}
    \left[
    \widehat{\dot{x}}_i(t)
    -
    F_{\theta}(x_i(t))
    -
    \sum_{j:(i,j)\in E} A_{ij}
    G_{\theta}(x_i(t),x_j(t))
    \right]^2.
\end{equation}
For the proposed HR experiments, $S$ is a degree stratum.  The intended compute
is controlled by the number of incoming edges to the stratum,
\begin{equation}
    M = |E_S| = |\{(i,j)\in E:i\in S\}|.
\end{equation}
Thus the scoring cost is $O(TMd)$ instead of $O(T|E|d)$.

\section{Implementation Issues Found and Corrected}

\begin{longtable}{p{0.04\linewidth} p{0.31\linewidth} p{0.34\linewidth} p{0.22\linewidth}}
\toprule
\# & Issue & Mathematical or experimental consequence & Correction \\
\midrule
\endfirsthead
\toprule
\# & Issue & Mathematical or experimental consequence & Correction \\
\midrule
\endhead
1 &
Candidate pool discarded optimized candidate parameters.  The worker trained
$(e,\theta)$, but the controller rebuilt fresh trees from $e$ before inserting
them into the pool. &
The score $S(e)$ was computed using one parameter state, while the stored
candidate used another.  Final fine-tuning started from random parameters
instead of the score-stage optimizer output. &
\verb|eval_candidate| now returns trained forcing and interaction FEX objects;
the pool stores those trained trees. \\
\addlinespace
2 &
Controller reward used $1/\sqrt{1+L}$ instead of $1/(1+L)$. &
This changes the reward curvature and therefore the policy-gradient advantage
used by the quantile controller. &
Reward changed to $S(e)=1/(1+L(e,\widehat{\theta}_e))$. \\
\addlinespace
3 &
No real stratified-controller scoring path existed. &
Previous stratified checks were fixed-design candidate scoring, not FEX
structure discovery.  They could not establish that the controller finds the
right operator sequence. &
Added \verb|score_nodes| through loss, FEX training, controller training, and
\verb|CoupledFEX.fit|. \\
\addlinespace
4 &
Stratified loss computed full-node tensors before slicing. &
Even when scoring only $S$, the code still evaluated self dynamics and
allocated interaction tensors over all $N$ nodes.  This undermined the claimed
$O(M)$ scoring cost. &
\verb|total_loss| now computes self dynamics only for selected nodes and
scatter-adds interactions into a tensor of size $|S|$. \\
\addlinespace
5 &
RBM required $N \bmod \texttt{num\_groups}=0$. &
This breaks ordinary settings such as $N=5000$ with batch size $p=32$, and it
is not required mathematically. &
RBM grouping now supports unequal group sizes. \\
\addlinespace
6 &
RBM inverse-probability weighting was initially treated like the default
baseline. &
The published RBM-CO baseline is raw RBM scoring followed by full-graph
fine-tuning.  IPW is a useful mechanism variant, but not the published method. &
Raw RBM is the default.  IPW is optional via \verb|rbm_reweight|. \\
\addlinespace
7 &
IPW needed exact probabilities for unequal batches. &
For unequal batches, using a single global $q$ is wrong.  The unbiased
correction depends on the retained batch size $s_b$. &
Optional IPW now uses $q_b=(s_b-1)/(N-1)$ per retained edge. \\
\addlinespace
8 &
Adam stored the best loss but left the model at the final epoch weights. &
The controller could reward a candidate for an earlier good epoch while
storing later worse parameters. &
Both network FEX and single FEX now clone and restore the best parameter state
before returning. \\
\addlinespace
9 &
LBFGS closure backpropagated the summed loss but returned the mean loss. &
Gradient scale did not match the reported objective, making Adam/BFGS scores
less comparable across batch counts. &
LBFGS closure now differentiates and returns the same mean loss. \\
\addlinespace
10 &
Network-FEX LBFGS did not use strong Wolfe line search, while single-FEX did. &
The two optimizer paths had different stability behavior. &
Network-FEX LBFGS now uses \verb|line_search_fn="strong_wolfe"|. \\
\addlinespace
11 &
Controller top-quantile pool could have size zero. &
If $\lfloor \rho K\rfloor=0$, the policy-gradient update has no candidates. &
The quantile pool size is now $\max(1,\lfloor \rho K\rfloor)$. \\
\addlinespace
12 &
GraphPool uniqueness check used the wrong node field. &
It checked \verb|node.op|, but nodes store operations as
\verb|node.operation.op|.  Uniqueness checking would fail when enabled. &
Fixed operation-name comparison in \verb|GraphPool.check_node_equality|. \\
\addlinespace
13 &
GraphPool checkpoint loading passed the wrong tree object. &
It passed \verb|tree_config.tree_func| into \verb|FEX|, but \verb|FEX| expects
the full \verb|TreeConfig|.  It also looked for sample indices in the wrong
dictionary level. &
Checkpoint loading now passes the full \verb|TreeConfig| and reads
\verb|_meta_sample_indices| from the saved state. \\
\addlinespace
14 &
\verb|CoupledFEX.to| used a nonexistent field. &
It referenced \verb|inter_fex|, but the candidate object stores
\verb|inter_tree|. &
Changed to \verb|inter_tree|. \\
\addlinespace
15 &
\verb|FEX.helpers| eagerly imported numerical-derivative code. &
The environment lacks \verb|pysindy|, so unrelated FEX imports could fail
because \verb|numerical_deriv| imported the smoothing dependency. &
\verb|numerical_deriv| is now lazy-loaded from \verb|FEX.helpers|. \\
\bottomrule
\end{longtable}

\section{Files Changed}

\begin{itemize}
    \item \verb|FEX/helpers/__init__.py|
    \item \verb|FEX/helpers/pools.py|
    \item \verb|FEX/training/loss_funcs.py|
    \item \verb|FEX/training/train_controller.py|
    \item \verb|FEX/training/train_fex.py|
    \item \verb|FEX/utils/fex.py|
    \item \verb|NumericalExperiments/HR/scripts/xingjian_test_experument.py|
\end{itemize}

\section{Verification Already Performed}

The following checks have passed:
\begin{itemize}
    \item Python compilation for all touched modules.
    \item A small candidate-evaluator smoke test returning trained FEX objects.
    \item A stratified-loss smoke test using \verb|score_nodes|.
    \item A non-divisible RBM grouping smoke test.
    \item \verb|git diff --check|.
\end{itemize}

These are smoke tests only.  They show that the corrected paths run, but they
do not yet establish recovery quality at the target network sizes.

\section{Remaining Gate Experiments}

\subsection{Controller Gate}

The first decisive experiment is a full controller run for HR at $N=5000$,
SNR 40:
\begin{equation}
    \text{stratified controller scoring}
    \quad\text{vs.}\quad
    \text{RBM-CO scoring}.
\end{equation}
The output should be reported in a Table-2 style format: each recovered
self-dynamics term $F$, each recovered interaction term $G$, its coefficient,
and final symbolic sMAPE.  This is the gate from fixed-design scoring to real
structure discovery.

\subsection{$k^\star$ U-Curve}

For degree stratum $k$, the predicted normalized risk is
\begin{equation}
    R(k)
    =
    \frac{g_0^2}{v}k
    +
    \frac{\sigma_\varepsilon^2}{v}\frac{1}{k}.
\end{equation}
The predicted optimum is
\begin{equation}
    k^\star
    =
    \frac{\sigma_\varepsilon}{|g_0|}.
\end{equation}
The important rule is to estimate and record
$\sigma_\varepsilon$, $g_0$, $v$, and $k^\star$ before sweeping degree strata.
If SNR decreases from $\infty$ to 45 to 40, then
$\sigma_\varepsilon$ should increase and $k^\star$ should move right.

\subsection{Assumption Checks}

Two checks bound the theorem:
\begin{enumerate}
    \item neighbor-state correlation on the attractor, which tests the
    conditionally i.i.d. neighbor-state assumption;
    \item low-degree state distribution versus global state distribution, which
    tests degree-state confounding in scale-free graphs.
\end{enumerate}

\subsection{Cost Panel}

The corrected cost panel should plot raw pair evaluations:
\begin{equation}
    C_{\mathrm{full}}\approx \bar{k}N,\qquad
    C_{\mathrm{RBM-CO}}\approx L N p,\qquad
    C_{\mathrm{strat}}\approx M.
\end{equation}
Here $L$ is the number of RBM stochastic partitions or controller score
realizations, $p$ is batch size, and $M$ is the fixed stratified edge budget.
This version makes RBM-CO climb above full sparse evaluation when $Lp>\bar{k}$.

\subsection{Moment-Matched Decoy}

The optional high-value test compares the true sigmoid coupling against a
variance-matched linear decoy.  The aggregation-collapse claim predicts that
hub scoring cannot separate these candidates, while leaf/low-degree scoring can.

\section{Current Status}

The implementation path is now suitable for launching the five tests, but the
paper claim is not yet established.  In particular, the full HR controller gate
at $N=5000$, SNR 40 has not yet been run.  The current result is therefore an
implementation correction and experimental harness preparation, not a final
scientific conclusion.

\end{document}
